\documentclass[11pt,oneside,openany]{book}
\usepackage[paperwidth=203.2mm,paperheight=254mm,top=20mm,bottom=22mm,left=21mm,right=19mm,headheight=15pt,headsep=7mm]{geometry}
\usepackage{fontspec}
\usepackage{xeCJK}
\setmainfont{texgyrepagella}[Extension=.otf,UprightFont=*-regular,BoldFont=*-bold,ItalicFont=*-italic,BoldItalicFont=*-bolditalic]
\setsansfont{texgyreheros}[Extension=.otf,UprightFont=*-regular,BoldFont=*-bold,ItalicFont=*-italic,BoldItalicFont=*-bolditalic]
\setmonofont{DejaVuSansMono.ttf}[Path=assets/fonts/,Scale=0.90]
\setCJKmainfont{NotoSerifCJKtc-Regular.otf}[Path=assets/fonts/,BoldFont=NotoSerifCJKtc-Bold.otf,ItalicFont=NotoSerifCJKtc-Regular.otf]
\setCJKsansfont{NotoSansCJKtc-Regular.otf}[Path=assets/fonts/,BoldFont=NotoSansCJKtc-Bold.otf]
\setCJKmonofont{NotoSansCJKtc-Regular.otf}[Path=assets/fonts/]
\usepackage{amsmath,amssymb,mathtools,cancel}
\usepackage{unicode-math}
\setmathfont{texgyrepagella-math.otf}
\let\setminus\smallsetminus
\renewcommand{\boldsymbol}[1]{\symbf{#1}}
\providecommand{\bm}[1]{\symbf{#1}}
\usepackage{newunicodechar}
\newunicodechar{∇}{\ensuremath{\nabla}}
\newunicodechar{∝}{\ensuremath{\propto}}
\newunicodechar{⋯}{\ensuremath{\cdots}}
\newunicodechar{✓}{\ensuremath{\checkmark}}
\newunicodechar{₀}{\textsubscript{0}}
\newunicodechar{₁}{\textsubscript{1}}
\newunicodechar{₂}{\textsubscript{2}}
\newunicodechar{ᵢ}{\textsubscript{i}}
\newunicodechar{⁻}{\ensuremath{^{-}}}
\usepackage{microtype}
\usepackage[table]{xcolor}
\definecolor{ink}{HTML}{162638}
\definecolor{accent}{HTML}{235F98}
\definecolor{soft}{HTML}{F1F4F7}
\usepackage{graphicx,adjustbox}
\usepackage{longtable,booktabs,array,calc,etoolbox}
\AtBeginEnvironment{longtable}{\footnotesize\setlength{\tabcolsep}{4pt}}
\usepackage{fvextra}
\fvset{fontsize=\footnotesize,breaklines=true,breakanywhere=true,breakautoindent=true,tabsize=4,frame=single,rulecolor=\color{black!15},framesep=5pt,baselinestretch=1}
\usepackage{enumitem}
\setlist{nosep,topsep=5pt,parsep=3pt,leftmargin=*}
\usepackage{titlesec}
\titleformat{\chapter}[display]{\sffamily\bfseries\color{ink}}{}{0pt}{\LARGE}
\titlespacing*{\chapter}{0pt}{0pt}{18pt}
\titleformat{\section}{\large\sffamily\bfseries\color{ink}}{}{0pt}{}
\titleformat{\subsection}{\normalsize\sffamily\bfseries\color{ink}}{}{0pt}{}
\titleformat{\subsubsection}{\normalsize\sffamily\bfseries}{}{0pt}{}
\setcounter{secnumdepth}{-1}
\setcounter{tocdepth}{0}
\usepackage{fancyhdr}
\pagestyle{fancy}\fancyhf{}
\fancyhead[L]{\small\sffamily\color{accent}Pixels → Perception}
\fancyhead[R]{\small\sffamily $edition$}
\fancyfoot[C]{\small\sffamily\thepage}
\renewcommand{\headrulewidth}{0.3pt}
\usepackage{xurl}
\usepackage[unicode,colorlinks=true,linkcolor=accent,urlcolor=accent,bookmarksdepth=3,pdftitle={Pixels to Perception — $edition$},pdfsubject={Website book edition},pdfcreator={LaTeX with Tectonic and Pandoc}]{hyperref}
\usepackage{bookmark}
\urlstyle{same}
\setlength{\parindent}{0pt}\setlength{\parskip}{5pt plus 1pt}
\setlength{\emergencystretch}{3em}
\raggedbottom
\providecommand{\tightlist}{\setlength{\itemsep}{0pt}\setlength{\parskip}{0pt}}
\providecommand{\real}[1]{#1}
\providecommand{\pandocbounded}[1]{#1}
\providecommand{\lt}{<}\providecommand{\gt}{>}
\newcommand{\bookchapter}[2]{\chapter*{#1}\phantomsection\label{#2}\addcontentsline{toc}{chapter}{#1}\normalfont\normalsize}
% A sub-chapter (course tree depth 2+): #3 is its contents level (section, subsection, …) and #4 that level's number; the
% contents show it (tocdepth is raised for this one line only) and the PDF bookmarks nest it under its main chapter.
\newcommand{\booksubchapter}[4]{\clearpage\phantomsection\label{#2}\addtocontents{toc}{\protect\setcounter{tocdepth}{#4}}\addcontentsline{toc}{#3}{#1}\addtocontents{toc}{\protect\setcounter{tocdepth}{0}}{\raggedright\Large\sffamily\bfseries\color{ink}#1\par}\vspace{14pt}\normalfont\normalsize}
\newcommand{\booksource}[2]{\par\begingroup\small\sffamily #1: \url{#2}\par\endgroup\smallskip}
\newcommand{\bookdisplay}[1]{\par\begingroup\centering\smallskip\begin{adjustbox}{max width=\linewidth}\(\displaystyle #1\)\end{adjustbox}\par\smallskip\endgroup}
\newcommand{\bookinline}[1]{$if(zh)$\allowbreak$else$\par\nobreak\noindent$endif$\begin{adjustbox}{max width=0.95\linewidth}\(#1\)\end{adjustbox}\allowbreak}
\newcommand{\bookimage}[1]{\par\begingroup\centering\smallskip\includegraphics[width=\linewidth,height=0.47\textheight,keepaspectratio]{#1}\par\smallskip\endgroup\nopagebreak[4]}
\renewenvironment{quote}{\begin{list}{}{\setlength{\leftmargin}{12pt}\setlength{\rightmargin}{8pt}}\item\small\color{ink}}{\end{list}}
$if(zh)$
\renewcommand{\contentsname}{目錄}
\renewcommand{\partname}{篇}
$endif$
\begin{document}
\hypersetup{pageanchor=false}
\begin{titlepage}
\color{ink}
\vspace*{8mm}
{\sffamily\small\bfseries WEBSITE BOOK EDITION\par}
\vspace{16mm}
{\color{accent}\rule{22mm}{4mm}\hspace{2mm}\rule{8mm}{4mm}\hspace{2mm}\rule{4mm}{4mm}\par}
\vspace{10mm}
{\fontsize{38}{43}\selectfont\bfseries Pixels\par → Perception\par}
\vspace{12mm}
{\sffamily\Large $subtitle$\par}
\vspace{5mm}
{\sffamily $tagline$\par}
\vspace{14mm}
{\sffamily\bfseries\Large $if(zh)$第 $volumeNumber$ 冊（共 $volumeCount$ 冊）$else$Volume $volumeNumber$ of $volumeCount$$endif$\par}
\vspace{2mm}
{\sffamily\large\color{accent} $volumeTitle$\par}
\vfill
{\sffamily\large $edition$\par}
\vspace{3mm}
{\sffamily $date$\par}
\vspace{8mm}
\rule{\linewidth}{0.5pt}
{\small\sffamily\url{https://blog.shyandram.dev/}\par}
\end{titlepage}
\frontmatter
\hypersetup{pageanchor=true}
\chapter*{$if(zh)$本書說明$else$About this edition$endif$}
$if(zh)$
本冊是 Pixels → Perception 網站完整課程分成的第 $volumeNumber$ 冊（共 $volumeCount$ 冊），主題為「$volumeTitle$」，收錄 $tutorials$ 篇已發布教學$if(posts)$與 $posts$ 篇部落格文章$endif$，依網站的學習單元順序編排。這是網站內容的書籍排版版本，並非另經同儕審查的出版品。

正文、程式碼、公式與參考文獻沿用網站內容。練習提示在本書中直接展開；靜態影像保留，互動示範則提供網站連結。原本的影像比較器改列為輸入與輸出影像。篇末參考文獻的編號僅適用於該篇教學。

本站內容仍在審閱中。製作過程使用 AI 協助，內容可能仍有錯誤或疏漏；請搭配原始文獻閱讀，並以網站的後續修訂為準。

兩種語言、各冊分別提供 PDF。連到同冊教學的連結會在本 PDF 內跳轉；連到其他冊的教學、特定小節、互動操作、筆記本與外部資料仍連至網站。
$else$
This is volume $volumeNumber$ of $volumeCount$ of the Pixels → Perception course, \emph{$volumeTitle$}. It contains $tutorials$ published tutorials$if(posts)$ and $posts$ blog posts$endif$, in the order of the website's learning tracks. This is a typeset edition of the website, not a separately peer-reviewed publication.

The prose, code, equations, and references retain the website's content. Exercise hints are expanded for print. Static figures remain in the book, image comparisons show their input and output images, and interactive demonstrations link to the website. Reference numbers are local to each tutorial.

The website is still under review. AI tools assisted its production, and errors or omissions may remain. Read it alongside the original sources and consult the website for subsequent corrections.

English and Traditional Chinese, and each volume, are separate PDFs. Links to tutorials in this volume navigate within the PDF. Links to tutorials in other volumes, to specific web sections, interactive demonstrations, notebooks, and external sources remain online.
$endif$
\bigskip
{\small\sffamily $if(zh)$內容快照$else$Content snapshot$endif$: $date$\par
$if(zh)$來源版本$else$Source revision$endif$: \texttt{$revision$}\par
\url{https://blog.shyandram.dev/}}
\tableofcontents
\mainmatter
$body$
\end{document}
